% Generated draft: requires mathematical and source review.
\documentclass[11pt]{article}
\usepackage[margin=25mm]{geometry}
\usepackage{fontspec}
\setmainfont{Latin Modern Roman}
\usepackage{amsmath,amssymb,mathtools,unicode-math}
\setmathfont{Latin Modern Math}
\usepackage{enumitem,needspace,xurl,hyperref}
\setlength{\parindent}{0pt}
\setlength{\parskip}{5pt}
\setlength{\emergencystretch}{3em}
\newcommand{\N}{\mathbb{N}}
\newcommand{\Z}{\mathbb{Z}}
\newcommand{\Q}{\mathbb{Q}}
\newcommand{\R}{\mathbb{R}}
\newcommand{\C}{\mathbb{C}}
\newcommand{\F}{\mathbb{F}}
\newcommand{\A}{\mathbb{A}}
\newcommand{\PP}{\mathbb{P}}
\newcommand{\E}{\mathbb{E}}
\newcommand{\eps}{\varepsilon}
\newcommand{\dd}{\,\mathrm{d}}
\begin{document}
\section*{UnsolvedMath ID: 30001974}
\textbf{Source problem code:} catalogue entry
\subsection{Definitions and mathematical statement}
The supplied record does not specify mathematical definitions of learning, tacit mathematical knowledge, or explicit mathematical knowledge. It supplies no mathematical domain, hypotheses, quantifiers, or criterion for measuring what an investigation gains by taking learning as its point of departure.

The source question is:
\begin{quote}
What does taking the learning of mathematics as a point of departure add to the investigation of tacit and explicit mathematical knowledge?
\end{quote}

This question is preserved as a historical, pedagogical, or methodological discussion prompt. The supplied statement-verification note characterizes it in those terms, marks its mathematical statement as \texttt{unrecoverable}, and promotes no separate clean statement. Consequently, the supplied evidence does not determine a mathematical proposition to formalize. Turning the question into a theorem or conjecture would require additional definitions and an intended mathematical target that are absent from the record. The registry's \texttt{open} label is supplied classification, not verification that this prompt is an unresolved mathematical proposition.
\subsection{Short English statement}
The question asks what studying mathematical learning contributes to understanding tacit and explicit mathematical knowledge. The supplied record treats it as a methodological discussion prompt and provides insufficient information for a precise mathematical proposition.
\subsection{Sources}
\begin{itemize}
\item Supplied local UnsolvedMath registry record: permanent ID \(30001974\), source problem code \texttt{catalogue entry}, title \emph{Mathematical Learning as Evidence of Tacit Knowledge}, collection \emph{Oberwolfach Reports Open Problems}. The recorded proposed year is 2012; no individual proposer is identified.
\item Supplied report citation: \emph{Explicit Versus Tacit Knowledge in Mathematics} (2012), Oberwolfach Reports, pp.~131--210, DOI \texttt{10.4171/owr/2012/04}, \url{https://doi.org/10.4171/owr/2012/04}. This is also the supplied literature-source link labelled ``original Oberwolfach report.'' The background describes Oberwolfach Reports as documenting mathematical workshops and their open-problem discussions.
\item Literature assessment dated exactly \texttt{2026-08-21}: the supplied curation labels the record \texttt{open} and assesses the question as asking what is gained by studying learning as a historical source, with no determinate mathematical truth value or settled resolution. This is supplied dated curation, not an independently verified literature assessment. The background lists the original report as checked and explicitly calls for independent verification of time-sensitive status and cited claims.
\item The supplied statement-verification note characterizes the official history-of-mathematics workshop report's question as historical, pedagogical, or methodological rather than an unresolved mathematical proposition. This characterization and the report attribution have not been independently verified here.
\end{itemize}
\end{document}
